\section{Desarrollo de la implementación}

Esta sección describe cómo se construyó el simulador de tráfico unidimensional con paralelismo OpenMP, explicando las decisiones de diseño y sus consecuencias en rendimiento. La idea central fue mantener una única base de código para la lógica de simulación y variar únicamente el mecanismo de ejecución y el nivel de optimización, de modo que cualquier diferencia de rendimiento sea atribuible al paralelismo o a la estrategia de compilación, y no a cambios en el algoritmo.

\subsection{Modelo del autómata celular de tráfico}

El sistema modela una carretera circular de $N$ celdas. Cada celda $i \in \{0, 1, \ldots, N-1\}$ puede estar vacía ($R = 0$) o contener un carro ($R = 1$). En cada paso de tiempo, un carro avanza a la celda siguiente si está libre, o se queda donde está si está ocupada. La regla de actualización es:

\begin{equation}
    R(t+1,\, i) =
    \bigl[R(t,i) \wedge R(t,\,(i+1)\bmod N)\bigr]
    \;\vee\;
    \bigl[R(t,\,(i-1)\bmod N) \wedge \neg R(t,i)\bigr].
    \label{eq:regla}
\end{equation}

El primer término representa un carro que no puede avanzar porque la celda siguiente está ocupada. El segundo representa un carro de la celda anterior que avanza hacia una celda libre. La condición $\neg R(t,i)$ en el segundo término evita que dos carros ocupen la misma celda simultáneamente.

En la implementación base, cada celda se almacena como un entero de 32 bits (\texttt{int}). La regla anterior se traduce directamente a operaciones de bits:

\begin{equation}
    \texttt{next}[i] =
    (\textit{here} \;\&\; \textit{ahead})
    \;\mid\;
    (\textit{behind} \;\&\; \mathtt{\sim}\textit{here} \;\&\; 1),
\end{equation}

donde \texttt{\&\,1} garantiza que el resultado sea siempre 0 o 1, ya que el operador \texttt{\textasciitilde} sobre enteros de 32 bits con signo produce \texttt{0xFFFFFFFF} en lugar de 1, lo que corrompería el conteo de carros.

Las lecturas del paso $t$ provienen del arreglo \texttt{cells} y las escrituras del paso $t+1$ van al arreglo \texttt{next\_cells}. Como son arreglos distintos y cada posición nueva depende únicamente de valores anteriores, todas las iteraciones del bucle son completamente independientes entre sí.

\subsection{Velocidad media y estado estacionario}

La velocidad media en el paso $t$ es la fracción de carros que se desplazaron:

\begin{equation}
    v(t) = \frac{1}{N_c}
    \sum_{i=0}^{N-1}
    \mathbf{1}\bigl[R(t,i) = 1 \;\wedge\; R(t+1,i) = 0\bigr],
    \label{eq:velocidad}
\end{equation}

donde $N_c$ es el número total de carros, que se conserva: la regla~\eqref{eq:regla} nunca crea ni elimina carros. Esto se verifica explícitamente en los tests de corrección.

Partiendo de una distribución inicial aleatoria, el sistema pasa por una fase transitoria donde $v(t)$ fluctúa mientras los carros se reorganizan. Tras suficientes pasos, $v(t)$ se estabiliza en la velocidad asintótica $v_\infty$, que es el valor de interés. El benchmark reporta el promedio sobre la fase de medición:

\begin{equation}
    v_\infty \approx \bar{v} =
    \frac{1}{T_m}
    \sum_{t = T_w + 1}^{T_w + T_m} v(t),
    \label{eq:vasint}
\end{equation}

donde $T_w$ es el número de pasos de calentamiento y $T_m$ el número de pasos de medición.

\subsection{Criterio de convergencia del calentamiento}

La duración de la fase transitoria varía con la densidad y con $N$. Un número fijo de pasos de calentamiento puede ser insuficiente para densidades cercanas al punto de transición ($\rho \approx 0.5$), o innecesariamente largo para densidades extremas donde el sistema converge rápido.

Para detectar la convergencia de forma automática, se compara la velocidad media en dos ventanas consecutivas de $W$ pasos. El calentamiento termina cuando:

\begin{equation}
    \left|
        \bar{v}_{[t-W,\, t]}
        - \bar{v}_{[t-2W,\, t-W]}
    \right| < \varepsilon,
    \label{eq:criterio}
\end{equation}

con $W = 50$, $\varepsilon = 10^{-4}$ y un mínimo de 200 pasos antes de evaluar el criterio, para evitar falsos positivos al inicio. El parámetro \texttt{max\_warmup\_steps} actúa como límite máximo; si la condición~\eqref{eq:criterio} no se cumple antes de ese límite, el calentamiento termina de todas formas. En el script de benchmark este techo se fija como $\max(N / 10,\; 2000)$, lo que garantiza al menos el tiempo suficiente para que una perturbación recorra la carretera completa. El número de pasos ejecutados se reporta en \texttt{stderr} como diagnóstico.

\subsection{Arquitectura de datos}

La estructura central es \texttt{Road}, que contiene el estado completo de la simulación:

\begin{itemize}
    \item \texttt{cells}: arreglo del estado actual, solo para lectura dentro de un paso.
    \item \texttt{next\_cells}: arreglo del siguiente estado, solo para escritura.
    \item \texttt{length}: número de celdas $N$.
    \item \texttt{car\_count}: número de carros $N_c$, constante durante toda la simulación.
\end{itemize}

Ambos arreglos se reservan una sola vez en \texttt{road\_create} y no se realizan asignaciones de memoria dentro del bucle de simulación, lo que elimina ese ruido de las mediciones de tiempo. Al finalizar cada paso, los punteros \texttt{cells} y \texttt{next\_cells} se intercambian directamente sin copiar datos.

\begin{figure}[H]
\centering
\fbox{%
\begin{minipage}{0.92\textwidth}
\centering
\textbf{Flujo de un paso de simulación}\\[4pt]
Leer \texttt{cells}
$\rightarrow$ Escribir \texttt{next\_cells} (regla~\eqref{eq:regla})
$\rightarrow$ Contar carros desplazados (ecuación~\eqref{eq:velocidad})
$\rightarrow$ Intercambiar punteros\\[2pt]
Si fase de calentamiento: evaluar criterio~\eqref{eq:criterio};
si no: acumular $v(t)$
\end{minipage}}
\caption{Ciclo compartido por todas las implementaciones.}
\label{fig:flujo_paso}
\end{figure}

El estado inicial se genera con \texttt{rand} usando como semilla \texttt{time(NULL)}. No se controla la semilla explícitamente porque la velocidad asintótica $v_\infty$ es estadísticamente estable respecto a la distribución inicial, como confirman las varianzas de \texttt{avg\_velocity} entre repeticiones en el CSV de resultados.

\subsection{Implementación secuencial}

El binario \texttt{traffic\_seq} se compila sin \texttt{-fopenmp}. Las directivas \texttt{\#pragma omp} presentes en \texttt{simulation.c} son ignoradas por el compilador, el binario no crea hilos en tiempo de ejecución y el bucle de actualización se ejecuta de forma ordinaria. Esto puede verificarse con \texttt{ldd bin/traffic\_seq}, que no lista \texttt{libgomp.so}.

La medición de tiempo usa \texttt{clock\_gettime(CLOCK\_MONOTONIC)}, un reloj que no se ve afectado por ajustes del sistema como NTP y garantiza mediciones fiables del tiempo transcurrido. El cronómetro se activa justo antes de la fase de medición y se detiene al terminar, dejando el calentamiento fuera del tiempo reportado.

La interfaz de línea de comandos es:

\begin{center}
\texttt{traffic\_seq} $N$ $\rho$ \texttt{max\_warmup} $T_m$
\end{center}

La salida contiene exclusivamente dos valores: el tiempo en milisegundos y la velocidad media $\bar{v}$. Esta separación permite al script de benchmark detectar ejecuciones que no alcanzaron el estado estacionario: una alta varianza en $\bar{v}$ entre repeticiones indica que el techo de calentamiento fue insuficiente.

\subsection{Implementación paralela con OpenMP}

El binario \texttt{traffic\_omp} se compila con \texttt{-fopenmp}. El número de hilos se lee como argumento de línea de comandos y se comunica al runtime con \texttt{omp\_set\_num\_threads}, lo que permite usar un único binario para todos los experimentos sin recompilar.

La interfaz extiende la del binario secuencial con un argumento adicional:

\begin{center}
\texttt{traffic\_omp} $N$ $\rho$ \texttt{max\_warmup} $T_m$ $p$
\end{center}

donde $p$ es el número de hilos. El valor efectivo se confirma con \texttt{omp\_get\_max\_threads()} y se registra en \texttt{stderr}.

Por paso de tiempo hay dos regiones paralelas. La primera, en \texttt{compute\_next\_generation}, aplica la regla~\eqref{eq:regla} a cada celda:

\begin{center}
\begin{verbatim}
#pragma omp parallel for schedule(static)
for (int i = 0; i < length; i++) { ... }
\end{verbatim}
\end{center}

La cláusula \texttt{schedule(static)} divide las celdas en bloques contiguos de tamaño fijo, uno por hilo. Esto es adecuado porque el trabajo por celda es uniforme, exactamente tres lecturas, dos operaciones lógicas y una escritura, y produce un patrón de acceso a memoria favorable para el procesador.

La independencia entre iteraciones es completa: todas las lecturas vienen de \texttt{current[]} y todas las escrituras van a \texttt{next[]}, que son arreglos distintos. No hay condiciones de carrera.

La segunda región paralela, en \texttt{count\_departed\_cars}, cuenta los carros desplazados:

\begin{center}
\begin{verbatim}
#pragma omp parallel for reduction(+:departed) schedule(static)
for (int i = 0; i < length; i++)
    departed += (old_cells[i] == 1 && new_cells[i] == 0) ? 1 : 0;
\end{verbatim}
\end{center}

La cláusula \texttt{reduction(+:departed)} hace que cada hilo lleve su propio contador local y los sume al final, evitando que los hilos se pisen entre sí. Entre las dos regiones paralelas no se necesita sincronización adicional: OpenMP garantiza automáticamente que todos los hilos terminaron la primera región antes de que cualquiera comience la segunda.

\subsection{Optimización por compilador}

Para cuantificar el efecto del compilador, se construyen dos variantes adicionales, \texttt{traffic\_seq\_opt} y \texttt{traffic\_omp\_opt}, a partir de los mismos archivos fuente, variando únicamente los flags de compilación.

El bucle principal realiza tres cargas de memoria, dos operaciones lógicas y una escritura por celda, lo que lo convierte en un kernel limitado principalmente por el ancho de banda de memoria. Las flags seleccionadas actúan sobre este perfil:

\begin{itemize}
    \item \texttt{-O3}: activa vectorización automática, transformaciones de bucle e inlining agresivo.

    \item \texttt{-march=native}: emite instrucciones AVX2 o AVX-512 según el procesador. Con datos \texttt{int} de 32 bits, un registro AVX2 procesa 8 celdas a la vez; con \texttt{uint8\_t} procesa 32.

    \item \texttt{-funroll-loops}: reduce el overhead del contador del bucle; más útil cuando los datos caben en L1 ($N \leq 64\,000$).

    \item \texttt{-flto} con \texttt{-fno-semantic-interposition}: permite al compilador fusionar e inlinar funciones entre archivos fuente distintos, eliminando el costo de llamada en cada uno de los $T_m = 1000$ pasos de medición.

    \item \texttt{-falign-loops=32}: alinea el inicio de cada bucle a 32 bytes, evitando que las instrucciones AVX2 crucen fronteras de línea de caché.
\end{itemize}

Las versiones base se compilan con \texttt{-O0} para servir como línea de referencia sin ninguna optimización.

\subsection{Optimización de memoria}

El profiling reveló que el cuello de botella de la implementación base es la cantidad de
instrucciones escalares. El IPC medido es aproximadamente 3.7, cercano al límite del
procesador, lo que indica que el CPU ejecuta instrucciones eficientemente sin esperar por
datos: el prefetcher secuencial anticipa los accesos y los resuelve mayoritariamente en L1,
como confirma el perfil de \texttt{perf mem} con una tasa de aciertos en L1 superior al
56\%. El problema no es que falten datos, sino que cada instrucción hace muy poco trabajo
útil: con celdas de tipo \texttt{int} de 32 bits, un registro AVX2 de 256 bits contiene
solo 8 celdas, y el compilador no puede vectorizar eficientemente un bucle cuya lógica
opera sobre enteros de ese ancho cuando la semántica real es binaria.

Para explotar esta oportunidad se implementaron \texttt{traffic\_seq\_mem} y
\texttt{traffic\_omp\_mem}, que aplican tres optimizaciones de forma conjunta.

\subsubsection{MEM-1: Reducción del tipo de celda a \texttt{uint8\_t}}

La estructura \texttt{RoadOpt} usa \texttt{uint8\_t~*} en lugar de \texttt{int~*}.
La lógica es idéntica pero se agrega un cast explícito en el complemento:

\begin{equation}
    \texttt{next}[i] =
    (h \;\&\; a) \;\mid\;
    (b \;\&\; \mathtt{(uint8\_t)(\sim h)} \;\&\; \mathtt{1u}),
\end{equation}

para evitar que la extensión de signo corrompa el resultado al escribir en un campo de
8 bits. El efecto inmediato es que un registro AVX2 de 256 bits pasa a contener 32 celdas
en lugar de 8, cuadruplicando el trabajo útil por instrucción SIMD. Como consecuencia
secundaria, el working set se reduce por un factor de 4, lo que beneficia los tamaños
de problema donde los datos con \texttt{int} superaban la capacidad de L3, como se
resume en la Tabla~\ref{tab:ws}.

\begin{table}[H]
\centering
\caption{Working set con \texttt{int} vs \texttt{uint8\_t} para dos buffers en Ryzen 5 7600X.}
\label{tab:ws}
\begin{tabular}{r r r l l}
\hline
$N$ &
\texttt{int} (bytes) &
\texttt{uint8\_t} (bytes) &
Nivel \texttt{int} &
Nivel \texttt{uint8\_t} \\
\hline
8\,000       & 64\,KB   & 16\,KB  & L2   & L1 \\
64\,000      & 512\,KB  & 128\,KB & L2   & L2 \\
500\,000     & 4\,MB    & 1\,MB   & L3   & L2 \\
2\,000\,000  & 16\,MB   & 4\,MB   & L3   & L3 \\
5\,000\,000  & 40\,MB   & 10\,MB  & DRAM & L3 \\
10\,000\,000 & 80\,MB   & 20\,MB  & DRAM & L3 \\
\hline
\end{tabular}
\end{table}

Para los tamaños donde ambas representaciones residen en L3 o superiores, la ganancia
proviene casi exclusivamente de la densidad SIMD. Para $N \geq 5\,000\,000$, donde
\texttt{int} desborda L3 hacia DRAM, la reducción del working set aporta un factor
adicional al eliminar accesos a memoria principal.

\subsubsection{MEM-2: Fusión de pasadas y eliminación del módulo}

La implementación base recorre el arreglo dos veces por paso: una para calcular el
siguiente estado y otra para contar los carros desplazados. La solución es fusionar
ambas pasadas en una sola función \texttt{compute\_next\_and\_count}, que lee y escribe
cada celda una sola vez, reduciendo a la mitad el número de recorridos sobre los datos:

\begin{center}
\begin{verbatim}
#pragma omp parallel for schedule(static) reduction(+:departed)
for (int i = 1; i < length - 1; i++) {
    uint8_t b = current[i - 1];
    uint8_t h = current[i];
    uint8_t a = current[i + 1];
    uint8_t n = (h & a) | (b & (uint8_t)(~h) & 1u);
    next[i]   = n;
    departed += (h == 1u && n == 0u) ? 1 : 0;
}
\end{verbatim}
\end{center}

Las celdas de los extremos ($i = 0$ e $i = N-1$), que requieren índices circulares, se
calculan fuera del bucle paralelo. Esto elimina el operador módulo del camino crítico y
deja el bucle interno libre de bifurcaciones adicionales, condición necesaria para que
el compilador pueda vectorizarlo con instrucciones AVX2.

\subsubsection{MEM-3: Alineación de buffers a línea de caché}

Los arreglos se reservan con \texttt{posix\_memalign(64, \ldots)} en lugar de
\texttt{malloc}, garantizando que el inicio de cada arreglo coincida con el inicio de
una cacheline de 64 bytes. Esto tiene dos efectos: los accesos en los extremos del
rango de cada hilo no afectan la línea de caché de otro hilo, evitando false sharing, y
el prefetcher puede anticipar los accesos más eficientemente desde una dirección base
alineada.

\subsection{Matriz de variantes de compilación}

Las cuatro estrategias anteriores producen ocho binarios:

\begin{table}[H]
\centering
\caption{Binarios generados por \texttt{make all} y sus características.}
\label{tab:binarios}
\begin{tabular}{llccc}
\hline
Binario & \texttt{impl} en CSV & OMP & OPT flags & Mem opt \\
\hline
\texttt{traffic\_seq}           & \texttt{traffic\_seq}           & No  & No  & No \\
\texttt{traffic\_seq\_opt}      & \texttt{traffic\_seq\_opt}      & No  & Sí  & No \\
\texttt{traffic\_omp}           & \texttt{traffic\_omp}           & Sí  & No  & No \\
\texttt{traffic\_omp\_opt}      & \texttt{traffic\_omp\_opt}      & Sí  & Sí  & No \\
\texttt{traffic\_seq\_mem}      & \texttt{traffic\_seq\_mem}      & No  & No  & Sí \\
\texttt{traffic\_seq\_mem\_opt} & \texttt{traffic\_seq\_mem\_opt} & No  & Sí  & Sí \\
\texttt{traffic\_omp\_mem}      & \texttt{traffic\_omp\_mem}      & Sí  & No  & Sí \\
\texttt{traffic\_omp\_mem\_opt} & \texttt{traffic\_omp\_mem\_opt} & Sí  & Sí  & Sí \\
\hline
\end{tabular}
\end{table}

\subsection{Diseño del script de pruebas}

El script \texttt{bench\_traffic.sh} actúa como harness de medición: su única responsabilidad es invocar los binarios con los parámetros correctos, capturar su salida y guardarla en un CSV.

\subsection{Matriz de configuraciones}

El harness define dos experimentos que cubren las preguntas de rendimiento centrales del trabajo.

El experimento de escalado mide todas las combinaciones de los seis valores de $N$ con los seis conteos de hilos $\{0, 2, 4, 6, 8, 12\}$ a densidad fija $\rho = 0.50$, para los ocho binarios de la Tabla~\ref{tab:binarios}. El valor $p = 0$ indica los binarios secuenciales; cualquier $p > 0$ invoca los binarios OpenMP con ese número de hilos. Los valores de $N$ cubren cada nivel de la jerarquía de memoria con la representación base; la Tabla~\ref{tab:ws} muestra cómo \texttt{uint8\_t} desplaza varios tamaños hacia niveles superiores de caché.

\subsection{Techo de calentamiento adaptativo}

El techo de pasos de calentamiento que el harness pasa a cada binario se calcula como:

\begin{equation}
    T_{w,\max} = \max\!\left(\frac{N}{10},\; 2000\right).
    \label{eq:techo}
\end{equation}

El factor $N/10$ garantiza que el techo sea proporcional al tiempo que tarda una perturbación en recorrer la carretera, con un margen de seguridad de un orden de magnitud. El piso de 2000 evita techos demasiado bajos para valores pequeños de $N$. El binario puede terminar antes si el criterio~\eqref{eq:criterio} se cumple, lo que ocurre en la práctica para densidades extremas ($\rho \leq 0.1$ o $\rho \geq 0.9$) en pocas centenas de pasos.

\subsection{Formato del CSV de resultados}

El archivo \texttt{data.csv} tiene siete columnas:

\begin{center}
\begin{tabular}{ll}
\hline
Columna & Descripción \\
\hline
\texttt{impl}           & Binario usado (véase Tabla~\ref{tab:binarios}) \\
\texttt{threads}        & Número de hilos (0 para variantes secuenciales) \\
\texttt{road\_length}   & Tamaño de la carretera $N$ \\
\texttt{density}        & Densidad inicial de carros $\rho$ \\
\texttt{repetition}     & Índice de repetición (1 a 10) \\
\texttt{wall\_time\_ms} & Tiempo de la fase de medición en ms \\
\texttt{avg\_velocity}  & Velocidad media $\bar{v}$ sobre los $T_m$ pasos \\
\hline
\end{tabular}
\end{center}

